\section{Requirements}\label{sec:3}

The requirements are the root of the traceability thread: each is a \textbf{testable} statement
with a stable ID, written in \textbf{EARS}\footnote{\textbf{E}asy \textbf{A}pproach to \textbf{R}equirements
\textbf{S}yntax, from Rolls-Royce's airworthiness work: A.~Mavin, P.~Wilkinson, A.~Harwood and
M.~Novak, ``Easy Approach to Requirements Syntax (EARS)'', \emph{2009 17th IEEE International
Requirements Engineering Conference}, pp.~317--322,
\href{https://doi.org/10.1109/RE.2009.9}{doi:10.1109/RE.2009.9}; practitioner guide at
\url{https://alistairmavin.com/ears/}. Of its six patterns this report uses four: the three
preconditions below, and the ubiquitous form that takes none.} form (an optional \texttt{When} event, \texttt{While}
state or \texttt{If} unwanted-condition precondition, then \emph{the system shall} a response),
traced to a design element in Sections \ref{sec:4}–\ref{sec:9} and to a verification in Section \ref{sec:10}. The authoritative
repository's copy is \texttt{project/REQUIREMENTS.md}; this section adds each
requirement's verdict and the regime it holds in.

The \textbf{verdict column is the honest one}, and it is the most important column in this report.
The offline evaluation harness has now run, so all seven requirements carry the experimental
evidence their acceptance criteria name, and \textbf{four of the seven carry a failing clause}:
R-1, R-2, R-3 and R-7, which the column says in its first words rather than at the end of a
paragraph. Two of those four fail on the harder of two arms and pass on the one their criterion
names, and both arms are stated. R-4's evidence is now a container kill as well as an in-JVM one,
and R-5's is a live rebuild \emph{and} an integration test that executes its SQL.

\begin{longtable}[]{@{}
  >{\raggedright\arraybackslash}p{(\columnwidth - 10\tabcolsep) * \real{0.05}}
  >{\raggedright\arraybackslash}p{(\columnwidth - 10\tabcolsep) * \real{0.07}}
  >{\raggedright\arraybackslash}p{(\columnwidth - 10\tabcolsep) * \real{0.06}}
  >{\raggedright\arraybackslash}p{(\columnwidth - 10\tabcolsep) * \real{0.28}}
  >{\raggedright\arraybackslash}p{(\columnwidth - 10\tabcolsep) * \real{0.25}}
  >{\raggedright\arraybackslash}p{(\columnwidth - 10\tabcolsep) * \real{0.29}}@{}}
\toprule
\textbf{ID} & \textbf{Type} & \textbf{Pri.} & \textbf{Statement (EARS)} & \textbf{Acceptance
criterion, and how it is verified} & \textbf{Verdict, and the regime it holds in} \\
\midrule
\endhead
R-1 & Func. & Must & \emph{When} a damper stops responding to commands, the system shall confirm the fault. &
p95 confirmation ≤ 15 min after the first ignored command; injection harness, ≥ 20 runs &
\textbf{Met as specified; not met in closed loop} \\*
R-2 & Func. & Must & \emph{While} a fault is latched, the system shall keep the affected zone below 1000 ppm by
compensation. & Zone CO₂ ≤ 1000 ppm for ≥ 95 \% of occupied minutes during the fault; end-to-end
scenario test & \textbf{Met on the case it covers; not met at design occupancy} \\*
R-3 & Non-f. & Must & The monitor shall raise no more than one false alarm per simulated week. & ≤ 1
\texttt{FaultConfirmed} in a 7-day fault-free replay, scored on the \emph{worst} week; replay test
on a fixed seed & \textbf{Not met} \\*
R-4 & Non-f. & Must & \emph{While} readings are stale beyond three intervals, the monitor shall report unavailable
and issue no diagnosis. & State is \texttt{BLIND}, zero fault events published; fault injection,
kill \texttt{co2-sensor} & \textbf{Met, in-JVM and at container level} \\*
R-5 & Non-f. & Must & The system shall make every diagnosis reconstructable. & Inputs, command, residual series,
classification and timestamp recoverable for any \texttt{caseId}; rebuild a case from the journal &
\textbf{Built, demonstrated live against the deployed stack, and now covered by integration
tests} \\*
R-6 & Non-f. & Should & The system shall identify the zone count at which the monitor misses its interval. &
Sustained zone count where the p95 monitor cycle exceeds 60 s is reported; scaling experiment to
breaking point & \textbf{Met, with an extrapolated answer}. And the deployed ceiling is a
different and far smaller number \\*
R-7 & Non-f. & Must & \emph{While} a fault is latched in one room, the system shall keep control and diagnosis of
every other room unchanged. & Zero fault events and no control-mode change in the other five, their
CO₂ \emph{and their airflow} inside the healthy-run band; injection in one room, ≥ 3 rooms rotated
as the injected one & \textbf{Met on the drawn weekday; not met at design occupancy} \\
\bottomrule
\end{longtable}

\textbf{No requirement here is Regulatory, and that is a gap rather than an absence.} Five of the
seven are non-functional, which is the right shape for a system whose interesting claims are about
latency, false-alarm rate, auditability, scale and isolation. But R-2's 1000 ppm is exactly the kind
of number a regulation supplies rather than a designer, and this report asserts it without a source.
Grounding it in the Swedish indoor-air guidance it is plainly borrowed from would convert R-2's
threshold from an invented parameter into a cited one, and §\ref{sec:13.3}, which owns the fact that every
physical parameter here is invented, would have one fewer row.

\textbf{Every measured figure behind those verdicts is in §\ref{sec:11.2}}, one row per requirement, and
\textbf{Appendix \ref{sec:A}} carries the same seven rows with their evidence column in full. The two failing
pairs are the ones to read twice: R-1 and R-7 pass on the arm their own criterion names and fail on
the harder arm beside it, and both arms are stated wherever either appears.

The verdict column above is the short form. \textbf{Appendix \ref{sec:A} reprints the same seven rows with
their full evidence column}, naming the verification each verdict rests on; §\ref{sec:11.2} reports the
experiments themselves.

\subsection{The tension these requirements set up}\label{sec:3.1}

R-2 and R-7 pull against each other, and that is deliberate. R-2 needs a physical lever that
reaches a faulted room from outside it; every such lever risks disturbing the rooms R-7 says must
stay unchanged. With six independent fans there would be no tension and nothing to design: R-7
would hold because nothing couples the rooms, and R-2 would have no mechanism. The shared air
handler creates the tension, and resolving it is the supervisor's whole job (Section \ref{sec:7.4}). R-1 and
R-3 pull against each other in the same way, through a single number: lowering the CUSUM limit
\texttt{h} detects sooner and cries wolf more often. Section \ref{sec:11.1} reports that trade-off as
measured rather than asserted, and §\ref{sec:11.2} now reports a sharper version of the same tension that was
not visible when this was written: the binding knob turned out to be not \texttt{h} but what the
CUSUM's \emph{reference} is, and the setting of it that meets R-3 outright takes R-1's closed-loop
detection from 15 of 20 to 1 of 20 (ADR-0010).

\subsection{Gaps with no requirement ID}\label{sec:3.2}

Kept visible rather than discovered during marking. Each is real work with no ID attached, so
nothing in the table above constrains it and no test can name an ID for it.

\begin{itemize}
\tightlist
\item \textbf{Data quality.} R-4 covers \emph{stale} readings only, and only insofar as staleness
  drives \texttt{BLIND}. Nothing requires duplicate suppression, out-of-order handling, or defines
  what the \texttt{quality} field in \texttt{readings/\{zone\}/co2} must mean.
\item \textbf{The dashboard.} No requirement states what the 3D view must show, how fresh it must
  be, or what counts as the fault being visible. Its verification is ``screenshot evidence for the
  report'', which is a deliverable and not a criterion.
\item \textbf{R-2's bound.} That a damper stuck fully shut must produce a work order rather than an
  attempted compensation is recorded in ADR-0008 and asserted by a unit test, but it is not itself
  a requirement, so no acceptance criterion covers it.
\end{itemize}

\textbf{Traceability.} The five architecture viewpoints map onto this document so that nothing is
left unverified: \emph{context} (4.1), \emph{functional} (4.2–4.5), \emph{information} (5–6),
\emph{behavioural} (8), \emph{deployment} (9). Section \ref{sec:10} closes the loop from each requirement
back to a verification, and Section \ref{sec:11.6} lists every verification that has not yet run.\label{sec:3end}
